\documentclass[11pt]{article}

\usepackage[a4paper, margin=1in]{geometry}
\usepackage{amsmath, amssymb}
\usepackage{bm}
\usepackage{booktabs}

\newcommand{\ReLU}{\operatorname{ReLU}}
\newcommand{\R}{\mathbb{R}}

\title{HydroShear: Original Implementation and Our UV-Metric Implementation}
\author{}
\date{}

\begin{document}

\maketitle

\section{Notation}

Let $\mathbf{p}_i \in \R^3$ be the $i$-th marker point on the elastomer
surface, and let
\begin{equation}
    \mathbf{x}_i =
    \begin{bmatrix}
        x_i \\
        y_i
    \end{bmatrix}
    \in \R^2
\end{equation}
be its tactile image coordinate. HydroShear decomposes the marker displacement
field into dilation and shear:
\begin{equation}
    \mathbf{M}_t(\mathbf{x},{}^E\mathbf{X}^I_{0:t})
    =
    \mathbf{M}^d_t(\mathbf{x})
    +
    \mathbf{M}^s_t(\mathbf{x},{}^E\mathbf{X}^I_{0:t}).
    \label{eq:field-decomposition}
\end{equation}
The original paper writes the tactile output as a 2D marker-flow field. The
local HydroShear code, however, computes an intermediate 3D vector field and
then projects or slices it to obtain a 2D image-space marker flow.

\section{Original HydroShear Implementation}

\subsection{Original Dilation}

Let $\phi_I(\cdot)$ be the signed distance function (SDF) of the indenter.
A tactile marker/source point is in contact with the indenter if
$\phi_I(\mathbf{p}_k)<0$. Define
\begin{align}
    C_t
    &=
    \{k \mid \phi_I(\mathbf{p}_k)<0\}, \\
    h^d_k
    &=
    -\phi_I(\mathbf{p}_k).
\end{align}
At the paper-formula level, the dilation displacement is written on the 2D
tactile plane:
\begin{align}
    \mathbf{v}_{ik}
    &=
    \mathbf{x}_i-\mathbf{x}_k, \\
    \mathbf{M}^{d}_{i,t}
    &=
    c_d
    \sum_{k\in C_t}
    h^d_k
    \mathbf{v}_{ik}
    \exp\{-\lambda_d\|\mathbf{v}_{ik}\|_2^2\}.
    \label{eq:original-paper-dilation}
\end{align}

The local original implementation uses the same Gaussian idea but computes the
intermediate displacement in 3D:
\begin{align}
    \mathbf{r}_{ik}^{3D}
    &=
    \mathbf{p}_i-\mathbf{p}_k, \\
    \mathbf{M}^{d,3D}_{i,t}
    &=
    c_d
    \sum_{k\in C_t}
    h^d_k
    \mathbf{r}^{3D}_{ik}
    \exp\{-\lambda_d\|\mathbf{r}^{3D}_{ik}\|_2^2\}.
    \label{eq:original-code-dilation}
\end{align}

\subsection{Original Hydroelastic Force Tracker}

Let $\mathbf{o}_{j,t}\in\R^3$ be the $j$-th indenter surface sample in the
elastomer frame at time $t$. Let $\phi_E(\cdot)$ be the elastomer SDF. The
original tracker first estimates how much of the object-point motion occurs
inside the elastomer:
\begin{align}
    \alpha_{j,t}
    &=
    \frac{
        \ReLU(-\phi_E(\mathbf{o}_{j,t}))
        -
        \ReLU(-\phi_E(\mathbf{o}_{j,t-1}))
    }{
        \phi_E(\mathbf{o}_{j,t})-\phi_E(\mathbf{o}_{j,t-1})+\epsilon
    }, \\
    \mathbf{d}_{j,t}
    &=
    \alpha_{j,t}
    \left(
        \mathbf{o}_{j,t}-\mathbf{o}_{j,t-1}
    \right).
    \label{eq:original-displacement}
\end{align}
With contact normal $\mathbf{n}_j$, the in-contact displacement is split into
normal and tangential components:
\begin{align}
    d^n_{j,t}
    &=
    \mathbf{d}_{j,t}^{T}\mathbf{n}_j, \\
    \mathbf{d}^{\tau}_{j,t}
    &=
    \mathbf{d}_{j,t}
    -
    d^n_{j,t}\mathbf{n}_j.
    \label{eq:original-normal-tangent-split}
\end{align}
The accumulated normal and tangential quantities are then
\begin{align}
    f^n_{j,t}
    &=
    \tilde{f}^{n}_{j,t-1}
    +
    E A_j d^n_{j,t}, \\
    \mathbf{f}^{\tau}_{j,t}
    &=
    \tilde{\mathbf{f}}^{\tau}_{j,t-1}
    +
    K A_j \mathbf{d}^{\tau}_{j,t}.
    \label{eq:original-force-update}
\end{align}
Normal pulling is removed, and the tangential component is clipped by Coulomb
friction:
\begin{align}
    \bar{f}^{n}_{j,t}
    &=
    \ReLU(f^n_{j,t}), \\
    \bar{\mathbf{f}}^{\tau}_{j,t}
    &=
    \min\left(
        1,
        \frac{\mu\bar{f}^{n}_{j,t}}
        {\|\mathbf{f}^{\tau}_{j,t}\|_2+\epsilon}
    \right)
    \mathbf{f}^{\tau}_{j,t}.
    \label{eq:original-friction}
\end{align}
Finally, the force-like tracked vector is reset when the object sample is not
inside the elastomer:
\begin{equation}
    \tilde{\mathbf{f}}_{j,t}
    =
    H(-\phi_E(\mathbf{o}_{j,t}))
    \left(
        \bar{f}^{n}_{j,t}\mathbf{n}_j
        +
        \bar{\mathbf{f}}^{\tau}_{j,t}
    \right).
    \label{eq:original-tracked-force}
\end{equation}

\subsection{Original Projection Point}

The paper distinguishes the contact force $\tilde{\mathbf{f}}_{j,t}$ from a
projection displacement $\hat{\mathbf{f}}_{j,t}$. The projected elastomer-side
contact point is
\begin{equation}
    \hat{\mathbf{o}}_{j,t}
    =
    \mathbf{o}_{j,t}
    +
    \hat{\mathbf{f}}_{j,t}.
    \label{eq:original-projected-point}
\end{equation}
Under the simplifying assumption $E=K$, $\mu=\hat{\mu}$, and uniform area
$A_j=A$, the paper relates the force and projection displacement by
\begin{equation}
    \tilde{\mathbf{f}}_{j,t}
    =
    KA\hat{\mathbf{f}}_{j,t}.
    \label{eq:force-projection-relation}
\end{equation}
This distinction is important because a physical force should not be directly
added to a position. In code, the calibrated tracked vector is often used as a
displacement-like quantity for shifting the Gaussian center.

\subsection{Original Shear}

Let
\begin{equation}
    K_t
    =
    \{j \mid \phi_E(\mathbf{o}_{j,t})<0\}
\end{equation}
be the set of in-contact indenter surface samples. In the paper-level 2D
formula, shear is aggregated from the projected contact location
$\hat{\mathbf{o}}_{j,t}^{xy}$:
\begin{align}
    h^s_j
    &=
    -\phi_E(\mathbf{o}_{j,t}), \\
    \mathbf{w}_{ij}
    &=
    \mathbf{x}_i-\hat{\mathbf{o}}_{j,t}^{xy}, \\
    \mathbf{M}^{s}_{i,t}
    &=
    c_s
    \sum_{j\in K_t}
    h^s_j
    \left(
        -\tilde{\mathbf{f}}^{xy}_{j,t}
    \right)
    \exp\{-\lambda_s\|\mathbf{w}_{ij}\|_2^2\}.
    \label{eq:original-paper-shear}
\end{align}

The local original implementation can be viewed as a 3D version. Let
$\mathbf{b}_{j,t}$ denote the calibrated displacement-like tracked vector used
by the code. The Gaussian source center is
\begin{equation}
    \mathbf{a}_{j,t}
    =
    \mathbf{o}_{j,t}
    +
    \mathbf{b}_{j,t}.
    \label{eq:original-code-center}
\end{equation}
Then the implemented 3D shear aggregation is
\begin{equation}
    \mathbf{M}^{s,3D}_{i,t}
    =
    c_s
    \sum_{j\in K_t}
    h^s_j
    (-\mathbf{b}_{j,t})
    \exp\{
        -\lambda_s
        \|\mathbf{p}_i-\mathbf{a}_{j,t}\|_2^2
    \}.
    \label{eq:original-code-shear}
\end{equation}
The original local code then forms
\begin{equation}
    \mathbf{M}^{3D}_{i,t}
    =
    \mathbf{M}^{d,3D}_{i,t}
    +
    \mathbf{M}^{s,3D}_{i,t}.
    \label{eq:original-code-total}
\end{equation}

\section{Our Current Implementation}

\subsection{Raycaster Coordinates and Metric}

For the curved Revo fingertip, every marker or object sample is associated with
a raycaster grid coordinate
\begin{equation}
    \mathbf{u}_i =
    \begin{bmatrix}
        u_i \\
        v_i
    \end{bmatrix}
    \in \R^2,
\end{equation}
where $u$ is the grid-column coordinate and $v$ is the grid-row coordinate.
Let $\mathbf{g}_u,\mathbf{g}_v\in\R^3$ be the world-space metric vectors
corresponding to a one-pixel step along the raycaster column and row
directions. In the current implementation these are global per-frame metric
vectors estimated from the ray grid:
\begin{equation}
    J
    =
    \begin{bmatrix}
        \mathbf{g}_u & \mathbf{g}_v
    \end{bmatrix}
    \in\R^{3\times2},
    \qquad
    G
    =
    J^T J
    \in\R^{2\times2}.
    \label{eq:our-metric}
\end{equation}
For two ray-grid points $\mathbf{u}_a$ and $\mathbf{u}_b$,
\begin{align}
    \Delta\mathbf{u}_{ab}
    &=
    \mathbf{u}_a-\mathbf{u}_b, \\
    \mathbf{r}^{uv}_{ab}
    &=
    J\Delta\mathbf{u}_{ab}, \\
    d^2_{uv}(a,b)
    &=
    \|\mathbf{r}^{uv}_{ab}\|_2^2
    =
    \Delta\mathbf{u}_{ab}^{T}G\Delta\mathbf{u}_{ab}.
    \label{eq:our-uv-distance}
\end{align}
This is a raycaster UV metric distance. It is more surface-aware than directly
using the 3D chord distance $\|\mathbf{p}_i-\mathbf{p}_j\|_2$, but it is still
not a full geodesic shortest-path distance over the mesh.

\subsection{Our Dilation}

In the current implementation, contact/penetration for marker source points is
estimated from the ray-depth field. Let $h^{uv}_k$ be this depth-based
penetration estimate. Our dilation field is
\begin{equation}
    \mathbf{M}^{d,uv}_{i,t}
    =
    c_d
    \sum_{k\in C_t}
    h^{uv}_k
    \mathbf{r}^{uv}_{ik}
    \exp\{
        -\lambda_d d^2_{uv}(i,k)
    \}.
    \label{eq:our-dilation}
\end{equation}
Compared with Eq.~\eqref{eq:original-code-dilation}, the Gaussian vector and
distance are no longer computed by the direct 3D difference
$\mathbf{p}_i-\mathbf{p}_k$. They are computed through the raycaster UV metric
in Eq.~\eqref{eq:our-uv-distance}.

\subsection{Our Force Split}

The recursive tracker still produces a per-object-sample tracked vector
$\mathbf{f}_{j,t}\in\R^3$. Let $\mathbf{n}_j$ be the contact normal used for
splitting this vector. We compute
\begin{align}
    q_j
    &=
    \mathbf{f}_{j,t}^{T}\mathbf{n}_j, \\
    h^{uv}_j
    &=
    \max(q_j,0), \\
    \mathbf{f}^{\tau}_{j,t}
    &=
    \mathbf{f}_{j,t}
    -
    q_j\mathbf{n}_j.
    \label{eq:our-force-split}
\end{align}
Here the normal component $q_j$ is used as the pressure/height weight, while
the tangential component $\mathbf{f}^{\tau}_{j,t}$ is used to shift and
propagate the shear field.

\subsection{Our UV Shear Center}

The current implementation moves the shear Gaussian center in raycaster UV
space using only the tangential component:
\begin{align}
    \Delta\mathbf{u}^{\tau}_{j}
    &=
    J^+
    \mathbf{f}^{\tau}_{j,t}, \\
    \hat{\mathbf{u}}_{j,t}
    &=
    \mathbf{u}_{j,t}
    +
    \Delta\mathbf{u}^{\tau}_{j}.
    \label{eq:our-uv-center}
\end{align}
Here $J^+$ is the pseudoinverse of the raycaster metric basis. This maps a
world-space tangential vector back to a 2D UV displacement.

\subsection{Our Tangent-Plane Vector Projection}

Before the source shear vector from object sample $j$ contributes to marker
$i$, it is projected onto marker $i$'s tangent plane:
\begin{align}
    P_i
    &=
    I
    -
    \mathbf{n}_i\mathbf{n}_i^T, \\
    T_{j\rightarrow i}
    \left(
        \mathbf{f}^{\tau}_{j,t}
    \right)
    &\approx
    P_i\mathbf{f}^{\tau}_{j,t}.
    \label{eq:our-tangent-projection}
\end{align}
This is an approximate tangent-plane projection. It is not a full parallel
transport operator along a geodesic.

\subsection{Our Shear}

The final UV-metric shear field is
\begin{equation}
    \mathbf{M}^{s,uv}_{i,t}
    =
    c_s
    \sum_{j\in K_t}
    h^{uv}_j
    \left(
        -P_i\mathbf{f}^{\tau}_{j,t}
    \right)
    \exp\{
        -\lambda_s
        d^2_{uv}(i,\hat{\mathbf{u}}_{j,t})
    \}.
    \label{eq:our-shear}
\end{equation}
The current total 3D marker displacement is
\begin{equation}
    \mathbf{M}^{uv}_{i,t}
    =
    \mathbf{M}^{d,uv}_{i,t}
    +
    \mathbf{M}^{s,uv}_{i,t}
    \in\R^3.
    \label{eq:our-total}
\end{equation}

\subsection{Our Visualization Projection}

For visualization, the 3D displacement is projected onto the normalized
raycaster axes:
\begin{align}
    \Delta u_i
    &=
    s_{\mathrm{px}}
    \mathbf{M}^{uv}_{i,t}
    \cdot
    \hat{\mathbf{g}}_u, \\
    \Delta v_i
    &=
    s_{\mathrm{px}}
    \mathbf{M}^{uv}_{i,t}
    \cdot
    \hat{\mathbf{g}}_v,
    \label{eq:our-visualization}
\end{align}
where
\begin{equation}
    \hat{\mathbf{g}}_u
    =
    \frac{\mathbf{g}_u}{\|\mathbf{g}_u\|_2+\epsilon},
    \qquad
    \hat{\mathbf{g}}_v
    =
    \frac{\mathbf{g}_v}{\|\mathbf{g}_v\|_2+\epsilon},
\end{equation}
and $s_{\mathrm{px}}$ is the pixel-per-meter arrow scale.

\section{What We Changed}

\subsection{Change 1: Euclidean 3D Kernel to Raycaster UV-Metric Kernel}

The original local implementation spreads dilation and shear with a direct 3D
Euclidean distance:
\begin{equation}
    \|\mathbf{p}_i-\mathbf{p}_j\|_2^2.
\end{equation}
Our implementation replaces this with the raycaster UV metric:
\begin{equation}
    d^2_{uv}(i,j)
    =
    (\mathbf{u}_i-\mathbf{u}_j)^T
    G
    (\mathbf{u}_i-\mathbf{u}_j).
\end{equation}
This makes the Gaussian propagation follow the ray-grid parameterization of the
curved fingertip instead of the straight 3D chord between two points.

\subsection{Change 2: Tangential Shift of the Shear Center}

The original code-level implementation shifts the shear Gaussian center with
the full tracked vector:
\begin{equation}
    \mathbf{a}_{j,t}
    =
    \mathbf{o}_{j,t}
    +
    \mathbf{b}_{j,t}.
\end{equation}
Our implementation first removes the normal component:
\begin{equation}
    \mathbf{f}^{\tau}_{j,t}
    =
    \mathbf{f}_{j,t}
    -
    (\mathbf{f}_{j,t}^T\mathbf{n}_j)\mathbf{n}_j,
\end{equation}
and then shifts the center in UV space:
\begin{equation}
    \hat{\mathbf{u}}_{j,t}
    =
    \mathbf{u}_{j,t}
    +
    J^+
    \mathbf{f}^{\tau}_{j,t}.
\end{equation}
Therefore, pressing into the sensor mostly changes the pressure weight, while
tangential motion shifts the shear source center.

\subsection{Change 3: Direct Vector Sum to Marker Tangent-Plane Projection}

The original code-level shear contribution uses the same source vector for all
markers:
\begin{equation}
    -\mathbf{b}_{j,t}.
\end{equation}
Our implementation projects the source tangential vector onto each target
marker's tangent plane:
\begin{equation}
    -P_i\mathbf{f}^{\tau}_{j,t}
    =
    -
    (I-\mathbf{n}_i\mathbf{n}_i^T)
    \mathbf{f}^{\tau}_{j,t}.
\end{equation}
This makes the propagated vector compatible with each marker's local surface
tangent plane.

\subsection{Change 4: Raycaster-Axis Visualization}

The original implementation ultimately reads the tactile field in a 2D sensor
plane. Our curved-fingertip implementation first computes a 3D displacement on
the surface and then projects it to the raycaster image axes:
\begin{equation}
    (\Delta u_i,\Delta v_i)
    =
    s_{\mathrm{px}}
    \left(
        \mathbf{M}^{uv}_{i,t}\cdot\hat{\mathbf{g}}_u,
        \mathbf{M}^{uv}_{i,t}\cdot\hat{\mathbf{g}}_v
    \right).
\end{equation}

\subsection{Current Limitations}

The current UV metric is still an approximation. It uses global per-frame
metric vectors $\mathbf{g}_u$ and $\mathbf{g}_v$, so it is not a fully local
surface metric $G(\mathbf{u})$. It also does not compute a true mesh geodesic
shortest path. In addition, the RevoLab implementation estimates contact depth
from the ray/grid surface representation instead of querying a true volumetric
elastomer mesh SDF.

\section{Compact Comparison}

\begin{center}
\begin{tabular}{lll}
\toprule
Component & Original implementation & Our implementation \\
\midrule
Dilation distance
&
$\|\mathbf{p}_i-\mathbf{p}_k\|_2^2$
&
$d^2_{uv}(i,k)=\Delta\mathbf{u}^T G\Delta\mathbf{u}$ \\
Shear center
&
$\mathbf{o}_{j,t}+\mathbf{b}_{j,t}$
&
$\mathbf{u}_{j,t}+J^+\mathbf{f}^{\tau}_{j,t}$ \\
Normal component
&
Can enter the full tracked vector
&
Used as pressure weight $h^{uv}_j$ \\
Shear vector
&
$-\mathbf{b}_{j,t}$
&
$-P_i\mathbf{f}^{\tau}_{j,t}$ \\
Surface handling
&
Direct 3D Euclidean propagation
&
Raycaster UV-metric propagation \\
Visualization
&
Planar tactile/image coordinates
&
Projection onto raycaster axes \\
\bottomrule
\end{tabular}
\end{center}

\end{document}
